%Time-stamp: <seg nov 17 20:21:36 2025 yair@ufabcdix>
\pdfoutput=1
\documentclass[aspectratio=169]{beamer}
\usetheme{CambridgeUS}
\usecolortheme{beaver}
\usefonttheme{professionalfonts}
\usepackage[brazil]{babel}
\usepackage[utf8]{inputenc}
\usepackage{amsmath,amssymb,mathrsfs}

\usepackage{bm}
\usepackage{graphicx}
\usepackage{booktabs}
\usepackage{tikz}
%%\usepackage{tkz-berge}
\usepackage{tkz-graph}
\usetikzlibrary{arrows,chains,fit,shapes,automata, positioning}
%\usetikzlibrary{arrows,shapes.gates.logic.US,calc}
\tikzset{>=latex}
\usepackage{xifthen}
\usepackage{randomwalk}

\usepackage{hyperref}
\hypersetup{
    colorlinks=true,
    linkcolor=blue,      % links internos
    urlcolor=blue,       % URLs
    citecolor=blue,      % citações
    filecolor=blue       % links de arquivos
}
\title{Algoritmo PageRank}

\author[MCBM022-23]{MCBM022-23 Introdução aos Processos Estocásticos }
\institute{UFABC}
\date{2025-3}

\begin{document}

%------------------------------------------------------------
\frame{\titlepage}

%------------------------------------------------------------
\begin{frame}{Motivação}
\begin{itemize}\itemsep=.75cm
  \item A \emph{Web} pode ser modelada como um \textbf{grafo dirigido}:
  \[
     G = (V,E), \quad V = \text{páginas}, \ E = \text{links entre páginas}.
  \]
  \item Problema: \textbf{como medir a importância de uma página?}
  \item Ideia do Google:
  \begin{quote}
  “Usar a estrutura de links da Web como um indicador de valor de uma página.”
  \end{quote}
  \item Cada link é interpretado como um “voto” de uma página a outra.
\end{itemize}
\end{frame}

%------------------------------------------------------------
\begin{frame}{O Modelo Probabilístico}
\begin{itemize}\itemsep=.5cm
  \item Considere o \textbf{surfista aleatório}: um usuário que clica em links ao acaso.
  \item Seja $G=(V,E)$ com $V=\{1,2,\dots,n\}$.
  \item Defina:
  \[
    \begin{aligned}
      N^+(a) &= \{b\in V : (a,b)\in E\}, \quad d^+(a)=|N^+(a)|\\
      N^-(a) &= \{b\in V : (b,a)\in E\}, % \quad d^+(a)=|N^+(a)|.
    \end{aligned}    
  \]
  \item O surfista move-se de $a$ para $b$ com probabilidade
  \[
    p_{a,b} =
    \begin{cases}
      \frac{1}{d^+(a)}, & \text{se }(a,b)\in E,\\
      0, & \text{caso contrário.}
    \end{cases}
  \]
\item Assim, $\mathsf{P}=(p_{a,b})$ é a matriz de transição de um
  passeio aleatório simples em $G$.
\end{itemize}
\end{frame}

%------------------------------------------------------------
\begin{frame}{Definição do PageRank}
\begin{itemize}
\item Cada página $a$ recebe uma pontuação $r_a$ proporcional à soma
  das pontuações das páginas que apontam para ela
  \[
  r_a = \sum_{b\in N^-(a)} \frac{r_b}{d^+(b)}.
  \]
  \item Em forma matricial \pause
  \[
     \bm{r} = \bm{r}\mathsf{P},
  \]
  isto é, $\bm{r} = (r_1,r_2,\dots,r_n)$ é um \textbf{vetor invariante}
  de $\mathsf{P}$.
\end{itemize}
\end{frame}

%------------------------------------------------------------
\begin{frame}{Problemas do modelo}
\begin{block}{1. Vértices pendentes (\emph{dangling nodes})}
  Páginas sem links de saída: $d^+(a)=0$.
  \begin{itemize}
    \item Fazem com que $\mathsf{P}$ não seja estocástica.
  \end{itemize}
\end{block}
\vfill
\begin{block}{2. Subconjuntos fechados (\emph{rank sinks})}
  \begin{itemize}
    \item Conjuntos de páginas que formam um bloco isolado.
    \item O surfista fica “preso” e nunca mais sai.
  \end{itemize}
\end{block}
\end{frame}

%------------------------------------------------------------
\begin{frame}{Exemplo de grafo \emph{Web}}
  \centering



      \begin{tikzpicture}[scale=1.32]
         \tikzstyle{vertex}=[circle,fill=black!25,minimum
         size=17pt,inner sep=0pt]
      
         \tikzstyle{edgeLabel} = [pos=0.5, text centered,
         font={\small}];
      
         \foreach \name/\pos/\text in {P-1/2/1, P-4/0/4}
         \node[vertex,xshift=6cm,yshift=.5cm] (\name) at (\pos,2)
         {$\text$};

         \foreach \name/\pos/\text in {P-3/2/3, P-5/0/5}
         \node[vertex,xshift=6cm,yshift=.5cm] (\name) at (\pos,0)
         {$\text$};

         \foreach \name/\pos/\text in {P-2/4/2, P-6/-2/6}
         \node[vertex,xshift=6cm,yshift=.5cm] (\name) at (\pos,1)
         {$\text$};

      
         \draw[->] (P-1) .. controls +(0:.7cm) and +(110:.7cm)
         .. (P-2) node[edgeLabel,above] {$1/2$};
      
         \draw[->] (P-1) .. controls +(-80:.7cm) and +(80:.7cm)
         .. (P-3) node[edgeLabel,right] {$1/2$};
      
         \draw[->] (P-2) .. controls +(170:.7cm) and +(-60:.7cm)
         .. (P-1) node[edgeLabel,above] {$1/2$};
      
         \draw[->] (P-2) .. controls +(-110:.7cm) and +(0:.7cm)
         .. (P-3) node[edgeLabel,below] {$1/2$};
      
         \draw[->] (P-3) .. controls +(135:1cm) and +(225:1cm)
         .. (P-1) node[edgeLabel,left] {$1/2$};
      
         \draw[->] (P-3) .. controls +(50:.7cm) and +(190:.7cm)
         .. (P-2) node[edgeLabel,below] {$1/2$};
      
         \draw[->] (P-4) .. controls +(0:1.2cm) and +(180:1.2cm)
         .. (P-1) node[edgeLabel,above] {$1/3$};

         \draw[->] (P-4) .. controls +(-80:1cm) and +(80:1cm) 
         .. (P-5)  node[edgeLabel,left] {$1/3$};

         \draw[->] (P-4) .. controls +(180:.7cm) and +(70:.7cm)
         .. (P-6) node[edgeLabel,above] {$1/3$};
      
         \draw[->] (P-5) .. controls +(0:1.2cm) and +(180:1.2cm)
         .. (P-3) node[edgeLabel,above] {$1/3$};

         \draw[->] (P-5) .. controls +(135:1cm) and +(225:1cm)
         .. (P-4) node[edgeLabel,left] {$1/3$};

         \draw[->] (P-5) .. controls +(180:.7cm) and +(-70:.7cm)
         .. (P-6) node[edgeLabel,below] {$1/3$};

         \draw (9,1.5) node[right] { \small $P=\begin{pmatrix}
              0&1/2&1/2&0&0&0\\
              1/2&0&1/2&0&0&0\\
              1/2&1/2&0&0&0&0\\
              1/3&0&0&0&1/3&1/3\\
              0&0&1/3&1/3&0&1/3\\
              0&0&0&0&0&0
           \end{pmatrix}$
         };
      \end{tikzpicture}


\vspace{1ex}
\small{Vértice 6 é pendente $\Rightarrow$ linha nula.}
\end{frame}

%------------------------------------------------------------
\begin{frame}{Correção dos vértices pendentes}
  \begin{itemize}\itemsep=.6cm
  \item Seja $\bm{u}\in\{0,1\}^n$ o vetor-linha indicador dos vértices
    pendentes e \(\mathbf{1}\) é o vetor-linha de 1's.
  \item Defina:
    \[
      \mathsf{Q} = \mathsf{P} + \bm{u}^\intercal \frac{1}{n}\mathbf{1}.
    \]
  \item Intuição: se o surfista cai numa página sem links, ele pula
    para uma página aleatória.
  \item Agora $\mathsf{Q}$ é uma matriz \textbf{estocástica}.
  \end{itemize}
\end{frame}

%------------------------------------------------------------
\begin{frame}{Matriz corrigida $\mathsf{Q}$}
\centering


      \begin{tikzpicture}[scale=1.3]
         \tikzstyle{vertex}=[circle,fill=black!25,minimum
         size=17pt,inner sep=0pt]
      
         \tikzstyle{edgeLabel} = [pos=0.5, text centered,
         font={\small}];
      
         \foreach \name/\pos/\text in {P-1/.2/1, P-4/-1.3/4}
         \node[vertex,xshift=6cm,yshift=.5cm] (\name) at (\pos,2)
         {$\text$};

         \foreach \name/\pos/\text in {P-3/.2/3, P-5/-1.3/5}
         \node[vertex,xshift=6cm,yshift=.5cm] (\name) at (\pos,0)
         {$\text$};

         \foreach \name/\pos/\text in {P-2/2.2/2, P-6/-3.3/6}
         \node[vertex,xshift=6cm,yshift=.5cm] (\name) at (\pos,1)
         {$\text$};

      
         \draw[->] (P-1) .. controls +(0:.7cm) and +(110:.7cm)
         .. (P-2) node[edgeLabel,above] {};
      
         \draw[->] (P-1) .. controls +(-80:.7cm) and +(80:.7cm)
         .. (P-3) node[edgeLabel,right] {};
      
         \draw[->] (P-2) .. controls +(160:.7cm) and +(-60:.7cm)
         .. (P-1) node[edgeLabel,above] {};
      
         \draw[->] (P-2) .. controls +(-110:.7cm) and +(0:.7cm)
         .. (P-3) node[edgeLabel,below] {};
      
         \draw[->] (P-3) .. controls +(100:1cm) and +(-100:1cm)
         .. (P-1) node[edgeLabel,left] {};
      
         \draw[->] (P-3) .. controls +(50:.7cm) and +(200:.7cm)
         .. (P-2) node[edgeLabel,below] {};
      
         \draw[->] (P-4) .. controls +(0:1.2cm) and +(180:1.2cm)
         .. (P-1) node[edgeLabel,above] {};

         \draw[->] (P-4) .. controls +(-80:1cm) and +(80:1cm) .. (P-5)
         node[edgeLabel,right] {};

         \draw[->] (P-4) .. controls +(180:.7cm) and +(70:.7cm)
         .. (P-6) node[edgeLabel,above] {};
      
         \draw[->] (P-5) .. controls +(0:1.2cm) and +(180:1.2cm)
         .. (P-3) node[edgeLabel,above] {};

         \draw[->] (P-5) .. controls +(100:1cm) and +(-100:1cm)
         .. (P-4) node[edgeLabel,left] {};

         \draw[->] (P-5) .. controls +(180:.7cm) and +(-70:.7cm)
         .. (P-6) node[edgeLabel,below] {};

         \draw[->,dashed] (P-6) .. controls +(0:.7cm) and +(180:.7cm)
         .. (P-2) node[edgeLabel,above] {$1/6$};

         \draw[->,dashed] (P-6) .. controls +(160:.7cm) and +(200:.7cm)
         .. (P-6) node[edgeLabel,above] {$1/6$};

         \draw[->,dashed] (P-6) .. controls +(0:.7cm) and +(210:.7cm)
         .. (P-1) node[edgeLabel,left] {$1/6$};

         \draw[->,dashed] (P-6) .. controls +(0:.7cm) and +(150:.7cm)
         .. (P-3) node[edgeLabel,left] {$1/6$};

         \draw[->,dashed] (P-6) .. controls +(0:.7cm) and +(210:.7cm)
         .. (P-4) node[edgeLabel,left] {$1/6$};

         \draw[->,dashed] (P-6) .. controls +(0:.7cm) and +(130:.7cm)
         .. (P-5) node[edgeLabel,left] {$1/6$};

         \draw (7.35,1.5) node[right] {\small  $Q=\begin{pmatrix}
              0&1/2&1/2&0&0&0\\
              1/2&0&1/2&0&0&0\\
              1/2&1/2&0&0&0&0\\
              1/3&0&0&0&1/3&1/3\\
              0&0&1/3&1/3&0&1/3\\
         1/6&1/6&1/6&1/6&1/6&1/6
      \end{pmatrix}$};
      \end{tikzpicture}

\vspace{1ex}
\small{O vértice 6 ``pula'' para qualquer outro uniformemente.}
\end{frame}

%------------------------------------------------------------
\begin{frame}{Tratamento dos subconjuntos fechados}
\begin{itemize}\itemsep=.5cm
  \item Ainda pode haver “sorvedouros” de probabilidade.
  \item Solução: permitir que o surfista reinicie ocasionalmente.
  \item Com probabilidade $
    \begin{cases}
      \alpha, &\text{ segue um link;}\\
      1-\alpha, &\text{ pula para uma página qualquer.}
    \end{cases}$
\end{itemize}
\[
\boxed{
  \mathsf{R} = \alpha\mathsf{Q} + (1-\alpha)\frac{1}{n}\mathbf{1}^\intercal\mathbf{1}
}
\]
\begin{itemize}
  \item $\mathsf{R}$ é positiva $\Rightarrow$ cadeia ergódica.
  \item Convergência garantida:
    $\bm{\sigma}^{(k)} \xrightarrow[k\to\infty]{} \mathbf{r}$, para
    qualquer distribuição inicial $\bm{\sigma}$.
\end{itemize}
\end{frame}

%------------------------------------------------------------
\begin{frame}{Tratamento dos subconjuntos fechados -- Exemplo de grafo \emph{Web}}
$$R=\begin{pmatrix}
  \frac{1-\alpha}{6} & \frac{\alpha}{2}+\frac{1-\alpha}{6} &
  \frac{\alpha}{2}+\frac{1-\alpha}{6} & \frac{1-\alpha}{6} &
  \frac{1-\alpha}{6} &
  \frac{1-\alpha}{6}\\[0.5em]
  \frac{\alpha}{2}+\frac{1-\alpha}{6} & \frac{1-\alpha}{6} &
  \frac{\alpha}{2}+\frac{1-\alpha}{6}
  & \frac{1-\alpha}{6} & \frac{1-\alpha}{6} & \frac{1-\alpha}{6}\\[0.5em]
  \frac{\alpha}{2}+\frac{1-\alpha}{6} &
  \frac{\alpha}{2}+\frac{1-\alpha}{6} & \frac{1-\alpha}{6}
  & \frac{1-\alpha}{6} & \frac{1-\alpha}{6} & \frac{1-\alpha}{6}\\[0.5em]
  \frac{1-\alpha}{6} & \frac{1-\alpha}{6} &
  \frac{\alpha}{3}+\frac{1-\alpha}{6} & \frac{1-\alpha}{6} &
  \frac{\alpha}{3}+\frac{1-\alpha}{6} & \frac{\alpha}{3}+\frac{1-\alpha}{6}\\[0.5em]
  \frac{1-\alpha}{6} & \frac{\alpha}{3}+\frac{1-\alpha}{6} &
  \frac{1-\alpha}{6} & \frac{\alpha}{3}+\frac{1-\alpha}{6} &
  \frac{1-\alpha}{6} & \frac{\alpha}{3}+\frac{1-\alpha}{6}\\[0.5em]
  \frac{\alpha}{6}+\frac{1-\alpha}{6} &
  \frac{\alpha}{6}+\frac{1-\alpha}{6} &
  \frac{\alpha}{6}+\frac{1-\alpha}{6} &
  \frac{\alpha}{6}+\frac{1-\alpha}{6} &
  \frac{\alpha}{6}+\frac{1-\alpha}{6} &
  \frac{\alpha}{6}+\frac{1-\alpha}{6}
\end{pmatrix}
$$

\end{frame}%------------------------------------------------------------
\begin{frame}{Tratamento dos subconjuntos fechados -- Exemplo de grafo \emph{Web}}
$$R=\begin{pmatrix}
0.025  & 0.45   & 0.45   & 0.025  & 0.025  & 0.025 \\[0.3em]
0.45   & 0.025  & 0.45   & 0.025  & 0.025  & 0.025 \\[0.3em]
0.45   & 0.45   & 0.025  & 0.025  & 0.025  & 0.025 \\[0.3em]
0.025  & 0.025  & 0.3083 & 0.025  & 0.3083 & 0.3083 \\[0.3em]
0.025  & 0.3083 & 0.025  & 0.3083 & 0.025  & 0.3083 \\[0.3em]
0.1667 & 0.1667 & 0.1667 & 0.1667 & 0.1667 & 0.1667
\end{pmatrix}
$$

\begin{center}
  $$\alpha = 0{,}85$$
\end{center}

\href{https://pagerank-ipe.streamlit.app/}{Resposta aqui}

\end{frame}
%------------------------------------------------------------
\begin{frame}{Cálculo iterativo do PageRank}
\begin{itemize}\itemsep=.65cm
  \item Método das potências:
  \[
    \bm{\sigma}^{(k+1)} = \bm{\sigma}^{(k)} \mathsf{R}.
  \]
  \item Convergência com taxa $|\lambda_2/\lambda_1|$.
  \item Nesse caso,  $\lambda_1=1$ e $ |\lambda_2| < \alpha$.
  \item Valor usado pelo Google: $\alpha = 0{,}85$.
\end{itemize}
\[
\Vert \bm{\sigma}^{(t+1)} - \bm{\sigma}^{(t)} \Vert < \varepsilon
\]
Critério de parada: $\varepsilon\in[10^{-2},10^{-7}]$, 50–100 iterações.
\end{frame}

%------------------------------------------------------------
\begin{frame}{Versão computacional}
  \begin{itemize}\itemsep=.5cm
  \item A matriz $\mathsf{P}$ é \textbf{esparsa}.
    
  \item É suficiente armazenar:  a estrutura de $\mathsf{P}$, o vetor $\bm{u}$
    
  \item A iteração prática:
    \[
      \bm{\sigma}^{(t)} = \alpha \bm{\sigma}^{(t-1)}\mathsf{P} + \alpha
      (\bm{\sigma}^{(t-1)}\mathbf{u}^\intercal)\frac{1}{n}\mathbf{1} +
      (1-\alpha)\frac{1}{n}\mathbf{1}
    \]onde \(\bm{\sigma}^{(t-1)}\mathbf{u}^\intercal\) é um escalar
    (massa atual em vértices pendentes) — assim evita-se multiplicar
    matrizes densas.
\end{itemize}
\end{frame}

%------------------------------------------------------------
\begin{frame}{Referências}
\small
\begin{itemize}
  \item Page, L.; Brin, S.; Motwani, R.; Winograd, T. \href{http://ilpubs.stanford.edu:8090/422/1/1999-66.pdf}{\textit{The PageRank Citation Ranking: Bringing Order to the Web.}} Stanford  (1998)
  \item Haveliwala, T. H. \href{https://nlp.stanford.edu/pubs/secondeigenvalue.pdf}{\textit{Second eigenvalue of the Google matrix.}} Stanford (2003)
\end{itemize}
\end{frame}

\end{document}
